
\documentclass[a4paper,12pt]{article}
\usepackage{amsmath}
\usepackage{amssymb}
\usepackage{geometry}
\geometry{margin=2.5cm}

\begin{document}

\section*{Electric Conduction – Solved Exercises}

\begin{enumerate}

% --------------------------------------------------------
\item 
Using the Richardson--Dushman equation
\[
J = A T^{2} \exp\!\left(-\frac{W}{k T}\right),
\]
the ratio of the emitted currents at temperatures $T_1 = 2200\,\text{K}$ and $T_2 = 2350\,\text{K}$ is
\[
\frac{J_2}{J_1}
= \left(\frac{T_2}{T_1}\right)^{2}
\exp\!\left[
- W \left(\frac{1}{kT_2} - \frac{1}{kT_1}\right)
\right].
\]
Substituting $W = 4.5\,\text{eV}$ and $k = 8.617 \times 10^{-5}\,\text{eV/K}$ yields
\[
\frac{J_2}{J_1} \approx 5.19.
\]

% --------------------------------------------------------
\item 
The retarding voltage reduces the emission current density as
\[
J = J_0 \exp\!\left(-\frac{|q|V}{kT}\right),
\]
and we require
\[
\frac{J}{J_0} = \frac{1}{15}.
\]
Thus
\[
V = \frac{kT}{|q|} \ln 15.
\]
For $T = 2200\,\text{K}$,
\[
V \approx 0.045\ \text{V}.
\]

% --------------------------------------------------------
\item 
The drift velocity in a copper wire is given by
\[
v_d = \frac{I}{n q A},
\]
where $A = \pi R^{2}$ with $R = 1.5\,\text{mm}$ and $n = 3.3 \times 10^{28}\,\text{m}^{-3}$.
Thus,
\[
v_d = \frac{12}{(3.3 \times 10^{28})(1.6 \times 10^{-19}) \pi (1.5 \times 10^{-3})^{2}}
\approx 2.04 \times 10^{-4}\,\text{m/s}.
\]

% --------------------------------------------------------
\item 
We examine an n-type semiconductor under the conditions:
\[
B_y = 0.4\,\text{T}, \qquad
U_H = 10\,\text{mV}, \qquad
I = 30\,\text{mA}.
\]

\begin{enumerate}
    \item 
    The Hall effect gives
    \[
    U_H = \frac{I B}{n q t},
    \]
    hence the electron concentration
    \[
    n = \frac{I B}{|q| U_H\, t}.
    \]
    (The thickness $t$ must be known to compute a numerical value.)

    \item 
    The drift velocity is
    \[
    v_d = \frac{I}{n q A},
    \]
    where $A$ is the cross-sectional area of the sample.

    \item 
    The Hall constant is
    \[
    R_H = \frac{1}{n q}.
    \]
\end{enumerate}

\end{enumerate}

\end{document}



\documentclass[a4paper,12pt]{article}
\usepackage{amsmath}
\usepackage{amssymb}

\begin{document}

\section*{Electric Conduction}

\begin{enumerate}
  \item By how much does the emitted current change if the tungsten cathode temperature is increased from $T = 2200\,\text{K}$ to $2350\,\text{K}$?  
  The work function for tungsten is $W_{\text{ki}} = 4.5\,\text{eV}$.  
  The Richardson coefficient is $6.02 \cdot 10^{5}\,\text{A}/\text{m}^2\text{K}^2$.

  \item What retarding voltage reduces the emission current density to one fifteenth for a cathode at $T = 2200\,\text{K}$?  
  The electron charge is $q = -1.6 \cdot 10^{-19}\,\text{As}$.

  \item What is the free-electron drift velocity in a copper conductor of radius $R = 1.5\,\text{mm}$ if the current is $I = 12\,\text{A}$?  
  The electron density in copper is $n_{\text{e}} = 3.3 \cdot 10^{28}\,\text{m}^{-3}$.

  \item We examine an n-type semiconductor under the conditions shown in the figure:  
  $B_{y} = 0.4\,\text{T}$, $U_{\text{H}} = 10\,\text{mV}$, $I = 30\,\text{mA}$.  
  The magnetic induction is directed along the $y$–axis.

  \begin{enumerate}
    \item What is the electron concentration?
    \item What is the drift velocity of the electrons?
    \item What is the Hall constant?
  \end{enumerate}

\end{enumerate}

\end{document}
